% ===================================================================
%  11. TAULA — Hiria eta bere eraikinak. — Sutea.
%  Traduction 2026 d'apres texto/io, controlee sur texto/fr, texto/en
%  et texto/es. Consigne : voir texto/eu/10-taula-01.tex.
%
%  OUVERTURE DE LA TROISIEME SERIE : « HIRUGARREN SAILA ».
%
%  LE BLOC c1-03-1 EST CELUI QUI A COUTE UNE INVERSION AU FINNOIS —
%  « la pordo (8) dil olda fortifikita kastelo (7) » — et il en coute
%  une au basque pour la meme raison, avec un remede different :
%  « atea (8), gaztelu gotortu zaharrarena (7) ».
% ===================================================================

%%K t11-apar-1 apar
\VUsaut{2.0mm}
\VUcentre{13.2pt}{90}{HIRUGARREN SAILA}
\VUsaut{3.0mm}
\VUfilet{16mm}
\VUsaut{5.6mm}
\VUtitre{15.0pt}{60}{11. TAULA}
\VUsaut{1.4mm}
\VUpk{11.4pt}{40}{Hiria eta bere eraikinak. --- Sutea.}
\VUsaut{2.2mm}

%%K t11-tit-1 sub
\VUpk{10.2pt}{40}{Inguruak.}
\VUsaut{3.0mm}

%%K t11-c1-01-1 p
1. --- Ozeanotik ehun kilometrora dagoen hiri handi batean bizi naiz, zeina hala
ere lehen mailako \VUgras{portua} baita \textsuperscript{(1)}. Haren ondotik doan
ibaiak laurehun metroko zabalera du eta kaialde bikaina osatzen.

%%K t11-c1-02-1 p
2. --- \VUgras{Kaien} \textsuperscript{(2)} ondoko hiriaren zati osoak erabat
itsas eta industria itxura du, eta \VUgras{aldiri} bat \textsuperscript{(3)}
osatzen, non marinelak eta langileak baino ez baitira bizi. Azken hauek
\VUgras{lantegietan} \textsuperscript{(4)} eta \VUgras{gas-fabrikan}
\textsuperscript{(5)} egiten dute lan, zeinaren \VUgras{tximinia} garaia
\textsuperscript{(6)} eta \VUgras{gasometroa} urrutitik ikusten baitira.

%%K t11-c1-03-1 p
3. --- Erdi Aroan hiria gotortuta zegoen, eta oraindik aurkitzen dira haren
harresien hondar batzuk, batez ere \VUgras{atea} \textsuperscript{(8)},
\VUgras{gaztelu gotortu zaharrarena} \textsuperscript{(7)}, bi
\VUgras{dorrerekin} \textsuperscript{(9)} alboetan, zeinak orain
\VUgras{espetxe} gisa erabiltzen baitira.

%%K t11-c1-04-1 p
4. --- \VUgras{Auzo} horretan guztian, oso zaharra dirudienean, eraikin bakar bat
da nahiko berria: \VUgras{aduana} \textsuperscript{(10)}. Kaiaren eta
\VUgras{kale} txiki baten \textsuperscript{(11)} artean dago, zeinak
\VUgras{udaletxe} zaharrera \textsuperscript{(12)} baitarama. Azken eraikin hori
erraz ezagutzen da bere \VUgras{zaindari-dorre} erraldoiagatik
\textsuperscript{(13)}, zeina hiri zaharraren gainetik altxatzen baita.

%%K t11-c1-05-1 p
5. --- Kalearen beste aldean aduanaren estilo bereko eraikin bat
\VUgras{dago}: \VUgras{burtsa} \textsuperscript{(14)} da. Haren fatxadetako bat
plaza txiki batera ematen du, non \VUgras{prefetura} baitago
\textsuperscript{(15)}. Gure hiria, egia esan, ez da hiriburua, baina hala ere
barrutiko gune garrantzitsua da.

%%K t11-tit-2 sub
\VUsaut{5.0mm}
\VUpk{10.2pt}{40}{Hiriaren goiko aldea.}
\VUsaut{3.4mm}

%%K t11-c1-06-1 p
6. --- Goarnizioa, nahiko handia, \VUgras{kuartel} zabal eta oso osasungarrietan
\textsuperscript{(16)} dago; \VUgras{hiriko parkeko} \textsuperscript{(17)}
burdinsareraino luzatzen dira. Parke horretara daraman
\VUgras{ibilbidea} \textsuperscript{(19)} \VUgras{aurrezki kutxaren}
\textsuperscript{(18)} atzetik igarotzen da eta \VUgras{katedralaren}
\textsuperscript{(20)} plazara ateratzen.

%%K t11-c1-07-1 p
7. --- Eraikin horiek guztiak mailaka altxatzen dira muino txiki batean, non gure
\VUgras{lizeo} handia ere \textsuperscript{(21)} baitago.

%%K t11-c1-07-2 p
Kale Nagusira ematen duen kale apur bat aldapatsutik jaisten bazara,
\VUgras{auzitegira} \textsuperscript{(22)} iritsiko zara, zeinaren aurrean
\VUgras{lorategitxo} atsegin bat \textsuperscript{(26)} baitago. \VUgras{Plaza}
hau ez da handia, baina hiriaren erdian berdetasun eta freskotasun oasi bat
osatzen du. Horregatik hiritarrek gogotik bisitatzen dute, eta atsegin dute
\VUgras{kafetegiaren} \textsuperscript{(25)} teilapean esertzea, ostegunetan eta
igandeetan \VUgras{kioskoan} \textsuperscript{(27)} jotzen duen musika-banda
entzuten, belardien eta lore-sailen artean.

%%K t11-c1-08-1 p
8. --- Belardi horietako batean brontzezko \VUgras{estatua} bat
\textsuperscript{(28)} dago, gure hiritarretako baten oroimenez jarritako
monumentua, zeinak zerbitzu handiak egin baitzizkion bere sorterriari.

%%K t11-tit-3 sub
\VUsaut{4.0mm}
\VUpk{10.2pt}{40}{Joan-etorria.}
\VUsaut{3.4mm}

%%K t11-c1-09-1 p
9. --- Gure hiria ez dago oraindik elektrizitatez argiztatua; \VUgras{gas-farolak}
\textsuperscript{(29)} baino ez ditu. Baina \VUgras{tokiko prentsa} argiztapen
modu hori hobetzeko lanean ari da, \VUgras{nola} jada hobetu baita tranbien
\VUgras{trakzioa}. \VUgras{Zaldiek} \VUgras{tiratzen zituzten} bagoi
\VUgras{zaharrak} \VUgras{tranbia elektriko} erosoez \textsuperscript{(31)}
ordezkatu dituzte, zeinek bizkor eramaten baitituzte bidaiariak
\VUgras{hiriaren mutur batetik bestera} oso prezio \VUgras{apalean}.

%%K t11-c1-10-1 p
10. --- Auzitegiaren ondoan \VUgras{bankua} dago \textsuperscript{(32)}, non
merkataritza-letrak deskontatzen baitira. \VUgras{Bankuko mandatariak}
\textsuperscript{(33)} etxez etxe joaten dira zor dena kobratzera.

%%K t11-tit-4 sub
\VUsaut{4.0mm}
\VUpk{10.2pt}{40}{Artea eta zientzia.}
\VUsaut{3.4mm}

%%K t11-c1-11-1 p
11. --- Hurbileko eraikin ederra \VUgras{museoa} da. Bertan daude batera hiriko
\VUgras{liburutegia} \textsuperscript{(34)}, \VUgras{natur zientzien bilduma}
bikainak \textsuperscript{(35)} (natur historiako museoa) eta \VUgras{pinakoteka}
dotorea \textsuperscript{(36)}, zeinak \VUgras{erakusketa} iraunkor bikaina
osatzen baitu.

%%K t11-c1-11-2 p
Astean bitan dago publikoarentzat zabalik, baina bisitariek edozein egunetan
ikus dezakete \VUgras{zaindariari} \textsuperscript{(37)} eskupeko txiki bat
emanda. \VUgras{Margolariak} \textsuperscript{(38)} lanera etortzen dira hara.
Kaballeteen \VUgras{aurrean} ikusten dituzu, \VUgras{paleta} eskuan, koadro
ospetsuak kopiatzen.

%%K t11-c1-12-1 p
12. --- Museoaren atzeko goietan ozta-ozta bereizten dira \VUgras{kupula}
\textsuperscript{(40)}, \VUgras{ospitalearena} \textsuperscript{(39)}, eta
\VUgras{irakasle-eskolaren} \textsuperscript{(41)} eraikinak, non gure hiriko
eskoletako irakasleek ikasi baitzuten. Antzoki-plazan \VUgras{posta}
\textsuperscript{(43)} dago, \VUgras{etxe partikular} batean
\textsuperscript{(42)} kokatua.

%%K t11-tit-5 sub
\VUsaut{5.0mm}
\VUpk{11.4pt}{40}{Sutea.}
\VUsaut{3.0mm}

%%K t11-tit-6 sub
\VUpk{10.2pt}{40}{«Sua!» oihua.}
\VUsaut{3.4mm}

%%K t11-c2-01-1 p
1. --- Hirian zehar berri kezkagarria zabaltzen da: antzokian sutea piztu berri
da! \VUgras{Plazara} \textsuperscript{(96)} iristen naizenean, jendetza bildu da
jada \VUgras{espaloian} \textsuperscript{(46)} eta \VUgras{galtzadan}
\textsuperscript{(47)}. \VUgras{Postako funtzionarioak} \textsuperscript{(44)}
eta \VUgras{postariak} \textsuperscript{(45)} postatik ateratzen ari dira.
\VUgras{Tranbia-gidariak} \textsuperscript{(48)} bere tranbia gelditu behar izan
du hiriko \VUgras{anbulantzia} \textsuperscript{(50)} igarotzen uzteko, zeina
\VUgras{kirurgialariarekin} \textsuperscript{(52)} eta \VUgras{erizainarekin}
\textsuperscript{(51)} etorri baita \VUgras{zaurituari} \textsuperscript{(53)}
laguntzera, zeina \VUgras{ohatilan} \textsuperscript{(54)} baitaramate
\VUgras{egunkari-kioskoaren} \textsuperscript{(55)} ondoko leku batera.

%%K t11-c2-02-1 p
2. --- Infanteria-konpainia bat, ariketatik itzultzen zena, gelditu egin da
ordena mantentzen laguntzeko \VUgras{hiriko zaindari} \textsuperscript{(61)}
zaldunekin batera. \VUgras{Zaldunek}, zeinen zaldiak hankaz gora altxatzen
baitira, \VUgras{soldaduek} \textsuperscript{(57)} eta \VUgras{jendarmeek}
\textsuperscript{(63)} baino hobeto moldatzen dira \VUgras{begiluzeak}
\textsuperscript{(56)} atzeratzen. Jendarmeak, gainera, oso lanpetuta daude.
Haietako batek \VUgras{poliziei} \textsuperscript{(64)} adierazten die nora eraman
beste \VUgras{kaltetua} \textsuperscript{(65)}, suhiltzaile ausart bat, eskaileratik
erori dena.

%%K t11-c2-03-1 p
3. --- \VUgras{Ofizialak} \textsuperscript{(66)} zehatz adierazten du non jarri
\VUgras{lurrun-ponpa} \textsuperscript{(67)}, zeina hurbiltzen ari baita. Makina
indartsu horren zain, \VUgras{esku-ponpa} \textsuperscript{(68)} jartzeko agindu
du, zeinaren \VUgras{mahuka} luzeak \textsuperscript{(69)} ur-uholdeak isurtzen
baititu sura. Sei suhiltzailek ponpatzen dute, eta haien lagunak,
\VUgras{turrusta-mutur} \textsuperscript{(70)} eusten, \VUgras{ur-zorrotada}
\textsuperscript{(71)} sutearen bihotz-bihotzera zuzentzen du.

%%K t11-c2-04-1 p
4. --- Antzokiaren beste aldetik \VUgras{eskailera} handi bat
\textsuperscript{(72)} altxatu dute. Suhiltzaileei sugarretara hurbiltzen uzten
die, zeinak \VUgras{ke} beltz moltsoen artean \textsuperscript{(75)}
altxatzen baitira.

%%K t11-tit-7 sub
\VUsaut{5.0mm}
\VUpk{10.2pt}{40}{Egintza ausartak.}
\VUsaut{3.4mm}

%%K t11-c2-05-1 p
5. --- \VUgras{Sua} \textsuperscript{(74)} \VUgras{antzokiaren}
\textsuperscript{(73)} atondoan hasi zen, \VUgras{opera} baten entseguan, zeinaren
lehen \VUgras{emanaldia} bihar egitekoa baitzen. Zorionez, aretoan
\VUgras{ikusle} gutxi batzuk \textsuperscript{(76)} baino ez zeuden. Jende gaixoa
\VUgras{balkoian} \textsuperscript{(77)} babestu da, nora \VUgras{suhiltzaile}
ausartak \textsuperscript{(86)} baitoazkie laguntzera eskaileraz eta sokaz.

%%K t11-c2-06-1 p
6. --- \VUgras{Arkupean} \textsuperscript{(78)}, non \VUgras{kartelak}
\textsuperscript{(79)} emanaldia iragartzen baitu, ikusten da nola doazen korrika
\VUgras{musikariak} \textsuperscript{(81)}, \VUgras{orkestrakoak}
\textsuperscript{(80)}, beren tresnak eramanez: \VUgras{biolina}
\textsuperscript{(81)}, \VUgras{txirula} \textsuperscript{(82)},
\VUgras{biolontxeloa} \textsuperscript{(83)}, \VUgras{tronboia}
\textsuperscript{(84)}, \VUgras{danborra} \textsuperscript{(85)} eta abar.
\VUgras{Altzarien} \textsuperscript{(87)} zati bat salbatzen saiatzen dira.

%%K t11-c2-07-1 p
7. --- \VUgras{Aktoreek} \textsuperscript{(89)} eta \VUgras{aktoresek}
\textsuperscript{(88)}, \VUgras{abeslariek} eta abeslarisek, eszenako
\VUgras{jantziekin} \textsuperscript{(90)} atera behar izan dute korrika. Tenorrak
\VUgras{kazetariari} \textsuperscript{(92)} azaltzen dio nola gertatu den.
\VUgras{Erreportaria} lehen unetik iritsi da agintariekin batera:
\VUgras{prefetarekin} \textsuperscript{(91)}, \VUgras{alkatearekin}
\textsuperscript{(93)}, \VUgras{polizia-buruarekin} \textsuperscript{(94)} eta
\VUgras{ikertzailearekin} \textsuperscript{(95)}, zeinak ikerketa egingo baitu
erantzukizuna noren gain dagoen zehazteko.
\VUsaut{6.0mm}
\VUfilet{22mm}
